\chapter{Mười Lát Cắt Nhìn Cho Công Bằng}
\markboth{Mười Lát Cắt Nhìn Cho Công Bằng}{Mười Lát Cắt Nhìn Cho Công Bằng}

\begin{truyen}{Mười Lát Cắt Rất Ngắn}{Ten Short Wake-Up Slices}
\chuhoa{P}{hần này không chửi điện thoại như đồ bỏ.}
\textit{(This section does not condemn the phone as pure trash.)}

Nó nhìn từng lát rất ngắn để thấy rõ một điều: điện thoại có lợi thật, nhưng lợi nhỏ không cứu nổi một cách sống dùng nó quá tay.
\textit{(It looks at short slices to show one thing clearly: the phone has real benefits, but small benefits cannot save a lifestyle that uses it without restraint.)}
\end{truyen}

\latcat{1. Gọi Cấp Cứu}{Điện thoại cứu được một mạng người khi cần gọi đúng lúc.\\
\textit{(A phone can save a life when it is used to call at the right moment.)}\\
Nhưng cũng chính cái máy đó có thể đẩy người ta vào tai nạn chỉ vì một tin nhắn vô nghĩa lúc đang lái xe.\\
\textit{(But that same device can also push someone into an accident because of a meaningless message while driving.)}}{A tool that saves in one minute can destroy in the next if its user has no discipline.}

\latcat{2. Từ Điển Trong Túi}{Tra từ vựng trong ba mươi giây là một ưu điểm rất thật.\\
\textit{(Looking up vocabulary in thirty seconds is a real advantage.)}\\
Nhưng nếu sau đó em trôi luôn thêm bốn mươi phút vào clip rác, cái lợi kia nhỏ quá để cứu nổi tổng thiệt hại.\\
\textit{(But if that is followed by forty minutes lost in junk clips, the benefit is too small to rescue the total damage.)}}{One clean drop does not purify a bucket of dirty water.}

\latcat{3. Bản Đồ Và Sự Lạc Đường}{Điện thoại giúp người ta tìm đúng đường ngoài phố.\\
\textit{(Phones help people find the correct road in the city.)}\\
Nhưng nhiều người đang ngày càng lạc đường trong chính nội tâm mình vì không còn nổi một khoảng yên để tự nghĩ.\\
\textit{(But many people are becoming more lost inside themselves because they can no longer endure even a quiet moment to think.)}}{Convenience outside does not guarantee clarity inside.}

\latcat{4. Lưu Giữ Kỷ Niệm}{Ảnh chụp giúp ta giữ lại một khoảnh khắc đẹp.\\
\textit{(Photos help preserve a beautiful moment.)}\\
Nhưng nếu cả đời chỉ lo ghi lại mà không sống sâu đủ để nhớ, thì mình đang tích file chứ chưa chắc đã tích đời.\\
\textit{(But if life becomes all recording and not enough deep living, we are collecting files rather than life itself.)}}{A memory kept by device is weaker than a moment truly lived.}

\latcat{5. Kết Nối Khẩn}{Một tin nhắn có thể nối người xa lại gần.\\
\textit{(A message can bring distant people closer.)}\\
Nhưng cũng có thể làm người ngồi cạnh mình biến mất ngay trước mắt vì mình không còn chịu ngẩng đầu lên.\\
\textit{(But it can also make the person beside us disappear because we no longer bother to look up.)}}{Fast connection can quietly murder real presence.}

\latcat{6. Làm Việc Nhanh}{Nhiều công việc nhờ điện thoại mà bớt chậm đi.\\
\textit{(Many tasks become faster thanks to phones.)}\\
Nhưng một đầu óc bị thông báo kéo liên tục thì làm nhanh ở chỗ nhỏ và suy yếu ở chỗ lớn.\\
\textit{(But a mind constantly dragged by notifications may become fast in small tasks and weak in larger thinking.)}}{Speed in fragments often reduces strength in depth.}

\latcat{7. Học Online}{Điện thoại mở ra bài giảng, tài liệu, ngoại ngữ và hàng ngàn lớp học.\\
\textit{(Phones open access to lessons, materials, languages, and thousands of classes.)}\\
Nhưng chúng cũng mở cùng lúc hàng vạn cánh cửa phân tâm, và phần lớn học sinh hiện nay không mạnh đến mức bước qua một cửa học mà bỏ qua chín mươi chín cửa còn lại.\\
\textit{(But they also open countless doors to distraction, and most students today are not strong enough to enter the learning door while ignoring the other ninety-nine.)}}{Access is not mastery. Open doors are not disciplined choices.}

\latcat{8. Tự Do Cá Nhân}{Điện thoại cho con người quyền chọn nội dung mình muốn xem.\\
\textit{(Phones give people the freedom to choose what they want to consume.)}\\
Nhưng thuật toán lại rất giỏi biến quyền chọn đó thành sự bị dắt mà vẫn tưởng là mình tự do.\\
\textit{(But algorithms are very good at turning that freedom into manipulation while still making people feel free.)}}{Not every choice we make is truly our own.}

\latcat{9. Giết Thời Gian}{Điện thoại làm lúc chờ đợi bớt chán.\\
\textit{(Phones make waiting feel less boring.)}\\
Nhưng chính việc không còn chịu nổi chán mới làm con người ngày càng nông và yếu sức chịu đựng.\\
\textit{(But the inability to endure boredom is exactly what makes people shallow and weak in endurance.)}}{The bored mind can grow. The constantly fed mind often only twitches.}

\latcat{10. Rất Hữu Ích, Nên Càng Nguy Hiểm}{Điện thoại không phải thứ vô dụng nên người ta mới lệ thuộc nó nhiều đến thế.\\
\textit{(Phones are not useless; that is precisely why people become so dependent on them.)}\\
Cái nguy của nó nằm ở chỗ: nó hữu ích vừa đủ để được tha thứ, và gây hại đủ nhiều để làm đời sống mòn đi từng ngày mà người ta vẫn tự bào chữa được.\\
\textit{(Its danger lies here: it is useful enough to be forgiven, and harmful enough to wear life down day by day while people still keep excusing it.)}}{A useful tyrant is harder to resist than a useless one.}

\ngancach